\section{Device Performance}

\begin{table}[htbp]
\centering
\caption{Photovoltaic parameters for control and SAM-treated devices
(average of 12 cells per condition; champion cell in parentheses).}
\label{tab:device-results}
\begin{tabular}{@{}lcccc@{}}
\toprule
\textbf{Device} & \textbf{PCE (\%)} & $V_{oc}$ \textbf{(V)} & $J_{sc}$ \textbf{(mA/cm}$^2$\textbf{)} & \textbf{FF (\%)} \\
\midrule
Control (no SAM) & 21.4 (21.9) & 1.088 & 24.8 & 79.1 \\
\textsc{sam-3} & 22.6 (23.0) & 1.121 & 24.9 & 80.9 \\
\textsc{sam-5} & 22.1 (22.5) & 1.104 & 24.7 & 80.6 \\
\textsc{sam-7} & 23.6 (24.1) & 1.150 & 25.0 & 82.3 \\
\bottomrule
\end{tabular}
\end{table}

Table~\ref{tab:device-results} summarizes photovoltaic parameters for the
control and three synthesized SAM candidates. \textsc{sam-7} --- a
phosphonic-acid-anchored molecule with a fluorinated terminal group ---
delivered the largest gains across all metrics, most notably a 62~mV
increase in $V_{oc}$ relative to the control, consistent with the
model's prediction that its combination of high binding energy and
moderate dipole moment would minimize interfacial recombination.
